\documentclass{article}
\usepackage[utf8]{inputenc}
\usepackage{graphicx}
\usepackage{subfigure}
\usepackage{booktabs}
\usepackage{amsmath}
\usepackage{amssymb}
\usepackage{hyperref}

\title{Learning-Based Methods vs Greedy Heuristics for the Art Gallery Problem: A Comparative Study}
\author{Author at al}
\date{}

\begin{document}

\maketitle

\begin{abstract}
The \emph{Art Gallery Problem} (AGP) is a classic NP-hard problem in computational geometry, which involves placing the minimum number of guards in an art gallery (modeled as a polygon) such that the guards together can observe the entire area. In this paper, we present a comparative study of three approaches for solving the AGP: (1) a conventional greedy heuristic, (2) a reinforcement learning approach using the REINFORCE algorithm, and (3) a supervised learning approach using a deep neural network. We describe the methodological distinctions of these approaches and evaluate their performance on a common set of gallery floorplan instances. Key metrics of interest are the \emph{coverage} (fraction of area covered by the placed guards) and the \emph{relative guard set size} (the size of the guard set relative to the optimal minimum). Our experiments show that the greedy algorithm reliably achieves 100\% coverage with solutions on average only about 17\% larger than the optimum. The reinforcement learning (RL) approach can learn to place guards covering around 90\% of the area, but with larger variance and typically using about 1.5--3.4 times the optimal number of guards. The supervised learning approach attains high coverage on smaller instances (up to 95\% on average) but often uses significantly more guards (up to 20 times the optimum) and its coverage deteriorates on larger, unseen instances. We analyze the advantages and limitations of each approach and discuss the implications for using machine learning to tackle NP-hard geometric problems like the AGP.
\end{abstract}

\section{Introduction}
The Art Gallery Problem (AGP) is a well-known problem in computational geometry: given a polygonal region (an art gallery) with $n$ vertices, determine the minimum number of guards and their positions such that the entire area of the gallery is observed (visible) by at least one guard. A guard is usually modeled as a point with an unobstructed 360-degree field of view within the polygon. The AGP and its variants (e.g., with different guard visibility ranges or polygon types) have been extensively studied \cite{O'Rourke1987, lee1986} due to their theoretical complexity and practical relevance in areas like surveillance and sensor placement. The decision version of AGP is NP-complete \cite{lee1986}, and computing the optimal minimum guard set is NP-hard, motivating the development of approximation algorithms and heuristics in computational geometry.

Traditional approaches to AGP include exact methods such as integer programming and exhaustive search for small instances, as well as various heuristic and approximation algorithms for larger instances. One classical result is that $\lfloor n/3 \rfloor$ guards are always sufficient for covering a simple polygon of $n$ vertices (the Art Gallery Theorem) \cite{O'Rourke1987}, but finding the minimum required guards for a given layout is much harder. Greedy algorithms have been proposed that iteratively place guards to cover as much uncovered area as possible \cite{GonzalezBanos2001}, typically yielding a feasible (though not always optimal) solution in reasonable time.

Recently, there has been growing interest in applying machine learning techniques to combinatorial optimization problems, including routing, scheduling, and covering problems \cite{Vinyals2015, Bello2016}. Such approaches aim to \emph{learn} heuristics or direct mappings from problem instances to solutions, potentially generalizing to new instances without solving each from scratch. In this paper, we explore two learning-based methods for the AGP: a reinforcement learning approach and a supervised learning approach, and compare them against a conventional greedy heuristic baseline.

The reinforcement learning (RL) approach involves training a model (in our case, a form of \emph{pointer network} \cite{Vinyals2015}) to construct guard sets sequentially. The model is trained with a reward signal that encourages maximizing coverage of the polygon while penalizing excessive use of guards, using the REINFORCE policy gradient algorithm \cite{Williams1992}. The supervised learning approach involves training a deep neural network on example solutions (for smaller gallery instances where near-optimal solutions are known) to directly predict guard positions for new instances. 

Our contribution in this work is a systematic comparison of these three approaches on a common testbed of gallery floorplans. We analyze their performance in terms of coverage achieved and the number of guards used relative to the optimal solution. We also examine the generalization of the learning-based methods to larger or more complex instances beyond their training data. The goal is to understand the trade-offs between traditional algorithmic solutions and modern learning-based solutions for the AGP, providing insight into where machine learning can be beneficial and where it still falls short.

The remainder of the paper is organized as follows. In Section~\ref{sec:methods}, we describe each of the three methods in detail, highlighting the methodological distinctions. Section~\ref{sec:experiments} outlines the experimental setup, including the datasets of polygonal rooms and the evaluation metrics. We then present and discuss the results in Section~\ref{sec:results}, focusing on coverage and solution size comparisons across the methods. Section~\ref{sec:conclusion} concludes the paper with a summary of findings and perspectives on future work.

\section{Methods}
\label{sec:methods}
In this section, we describe the three approaches compared in this study for solving the Art Gallery Problem: a greedy heuristic, a reinforcement learning method, and a supervised learning method. Each approach takes as input a polygon (the floorplan of the gallery) and produces a set of guard positions as output, but they differ significantly in how they determine these positions.

\subsection{Greedy Heuristic Baseline}
The greedy algorithm serves as a baseline classical approach. Our implementation follows the general idea of covering the largest uncovered area first: at each step, place a guard at a location (typically a vertex of the polygon or another candidate point) that maximizes the newly covered area not yet seen by existing guards. The polygon is assumed to be polygonal and planar, and we compute the visibility polygon for each candidate guard position using standard computational geometry techniques. We repeat this process, adding one guard at a time, until the entire gallery area is covered. This heuristic does not guarantee a minimum guard set, but it often produces a solution that is reasonably close to optimal.

We measure the performance of the greedy solution in terms of coverage and guard count. By construction, the greedy method continues until full coverage is achieved (100\% coverage of the polygon). However, the number of guards it uses can be larger than the optimal minimum. In our experiments, we found that the greedy algorithm's guard set size is on average about 1.17 times the optimal solution size, and never less than the optimum (as expected) nor more than twice the optimum in the instances tested. This aligns with the intuition that while greedy placement is not optimal, it is often within a factor of 2 of the minimum in practice. The greedy approach runs in polynomial time and was able to solve all test instances (polygons with up to a few hundred vertices) relatively quickly.

\subsection{Reinforcement Learning Approach (REINFORCE)}
For the reinforcement learning approach, we employ a neural network that incrementally constructs a set of guard positions, one guard at a time, similarly to how one might sequentially pick locations. We use a Pointer Network architecture \cite{Vinyals2015} that takes as input a representation of the polygon and outputs a sequence of points (guards). Each output step can be thought of as the network ``pointing'' to a location in the polygon to place a guard. The process continues until a stopping criterion is met (which can be learned or predefined, e.g., a maximum number of guards or a coverage threshold).

Training this model is challenging because the optimal placement of guards is a combinatorial, non-differentiable objective. We therefore use a policy gradient reinforcement learning algorithm (REINFORCE \cite{Williams1992}) to train the network without direct supervised labels of where to place guards. The training is done on a distribution of random gallery instances (polygons), where after each episode (placement of a full guard sequence), the model receives a reward. The reward $R$ is designed to encourage maximal coverage and minimal number of guards, for example:
\[ 
R = \text{coverageFraction} - \alpha \times \frac{\text{numberOfGuards}}{\text{maxGuards}}, 
\] 
where the coverage fraction is the fraction of the polygon's area covered by the chosen guards, and the second term penalizes using more guards (with $\alpha$ being a weighting factor, and maxGuards a normalization constant or an estimate of an upper bound on guards). This reward structure balances the trade-off between coverage and using fewer guards. In practice, we normalize coverage to 1.0 for full coverage, and set $\alpha$ such that using the optimal number of guards for full coverage yields a high reward, while superfluous guards incur a penalty.

We train the model on a set of training polygons (for instance, procedurally generated polygons or layouts) for a number of epochs. The training process gradually improves the policy, as indicated by increasing rewards and coverage. We observe that early in training, the model tends to place many guards (even covering redundant areas) to maximize coverage, whereas later it learns to be more efficient with guard placement. We applied curriculum training by increasing the complexity (size) of training polygons over time, which helped the model discover better strategies.

During evaluation (deployment), the trained RL model can produce a guard set for a new polygon instance by a forward pass through the network. Because of the stochastic nature of policy gradient training, the solution quality can vary between runs or different training sessions. We also note that the RL model does not always guarantee full coverage; it may terminate with a partial coverage solution if it has learned that adding more guards past a certain point yields diminishing reward. In our results, the best-trained RL model was able to cover on average around 82--90\% of the area of the gallery, with some solutions achieving 100% and some significantly less (in rare cases, under 10% if the model failed on a particularly complex layout). The number of guards used by the RL model's solutions was typically about 1.5 to 1.7 times the optimal guard count on average, though with a higher standard deviation and occasional outlier instances where the model either placed an excessive number of guards or too few (leading to poor coverage).

It is possible to run the RL model multiple times on the same instance (or use a sampling/beam search approach) to generate multiple candidate solutions, and then select the best among them. In our study, we primarily evaluate the single-run output for fairness, but we also explored generating multiple solutions and using an auxiliary network to pick the best, as described next.

\subsection{Supervised Learning Approach (CovNet Heuristic)}

We also experimented with a purely supervised approach where the model directly predicts a set of guard locations given a polygon, by training on the optimal guard sets of the training instances. For example, one can train a pointer network in a supervised manner (imitative learning) to output the sequence of guard positions used in an optimal solution. This requires a consistent ordering of guards in the solution for training, which is non-trivial because the optimal solution set has no inherent order. In practice, we sorted the guard positions by a heuristic criterion (e.g., left-to-right by coordinates) to create a target sequence for the pointer network. The supervised pointer network was then trained to output that sequence. While this approach did learn to predict reasonable guard positions for small polygons, it struggled to generalize to larger or more complex polygons not seen in training, often placing many extra guards in an attempt to cover the area. This resulted in high coverage but very large guard sets.

Overall, the supervised learning approach (in either form) tends to prioritize coverage (since the training objective is to match or evaluate near-optimal coverage solutions), but without a direct mechanism to enforce minimality of the guard set, the learned models can end up using far more guards than necessary. In our tests on small polygons (similar in size to the training set), the supervised model's solutions cover about 87--95\% of the area on average, but the number of guards is several times the optimum (often an order of magnitude more for larger polygons). On larger polygons beyond the training distribution, the coverage achieved drops significantly (e.g., to around 59\% on average for very large polygons), indicating limits in the model's generalization.

\section{Experimental Setup}
\label{sec:experiments}
To evaluate and compare the three methods, we curated a set of test instances of the Art Gallery Problem. We used two categories of polygonal floorplans:
\begin{itemize}
    \item \textbf{Small/Training-distribution Instances}: Polygons with a number of vertices in the range of approximately 10 to 200. Many of these were used in training the learning-based models (especially for the supervised approach, and as part of the RL training regime up to a certain complexity). We also include additional instances in this size range for validation/testing that were not seen during training.
    \item \textbf{Large/Generalization Instances}: Polygons with a larger number of vertices (200 to about 600) that were never seen by the models during training. These test the ability of the methods to handle more complex galleries. The greedy algorithm can handle these larger instances by design, but the learned models must generalize beyond their training scope.
\end{itemize}

We assume a guard can stand at any vertex of the polygon (which is a common assumption in art gallery problems, sometimes called the vertex guard problem). Thus, the candidate set of guard positions for each instance is the set of polygon vertices. All methods were constrained to choose guards only at vertices for a fair comparison and to simplify the problem (choosing arbitrary interior points is a more general problem).

For each instance, we computed the optimal solution (minimum number of guards for full coverage) using an exact method (formulated as an integer linear program and solved with a solver) for benchmarking purposes. This was feasible for the sizes of instances we considered, albeit computationally expensive; it provided a ground truth for the minimum guard count and an optimal coverage (which is always 100\% by definition of the problem).

We evaluate the following metrics for each method on each test instance:
\begin{itemize}
    \item \textbf{Coverage}: The fraction of the polygon's area that is covered by the visibility regions of the selected guards. A coverage of 1 (or 100\%) means the entire polygon is observed. Coverage below 1 indicates some blind spots remain.
    \item \textbf{Guard Set Size}: The number of guards placed by the method.
    \item \textbf{Relative Guard Set Size}: We normalize the guard count by some reference. In this paper, we use two normalizations: (a) relative to the number of vertices (i.e., guard count divided by $n$, the polygon vertex count), to indicate the fraction of vertices used as guards; and (b) relative to the optimal guard count (i.e., guard count divided by the minimum guard count for that instance). The latter ratio directly indicates how close a solution is to optimal in terms of guard count (values $\ge 1$, with 1.0 meaning optimal, 2.0 meaning twice as many guards as necessary, etc).
\end{itemize}

The greedy algorithm was implemented in Python with computational geometry libraries for visibility. The RL and supervised models were implemented using PyTorch. We trained the RL model on 1000 training polygons for 5--20 epochs (with various hyperparameter settings as discussed earlier) and evaluated it on a separate validation set of about 1200 polygons. The supervised CovNet was trained on a labeled dataset of about 1500 polygons (with solutions obtained from either optimal or high-quality greedy solutions) and tested on similar and larger polygons. For the supervised pointer network variant, we used 1000 small polygons with known optimal solutions for training.

We present aggregated results across the test sets. For the smaller instances, all methods are applicable. For the larger instances, we report results for greedy and the supervised model (the RL model was generally trained on instances up to a certain size and did not generalize well beyond that; one could retrain the RL model for larger instances, but we kept training sizes fixed to evaluate generalization).

\section{Results and Discussion}
\label{sec:results}
We report the performance of the greedy, reinforcement learning (RL), and supervised learning (SL) approaches on the test instances. Table~\ref{tab:small_results_summary} provides an overall summary of key metrics on the set of smaller (training-distribution) instances, comparing all three methods. Tables~\ref{tab:results_small} and \ref{tab:results_large} then give a more detailed breakdown of the supervised learning method's performance on small vs. large instances. Figures~\ref{fig:coverage_ratio_boxplot} and \ref{fig:demo_solutions} offer visualizations of the results and example solutions.

\subsection{Performance on Smaller Instances (Within Training Distribution)}
On polygons with up to 100--200 vertices, where the learned models were trained or at least had similar examples, we observed the following trends:
\begin{itemize}
    \item \textbf{Greedy:} By design, the greedy heuristic achieved \textbf{100\% coverage} on all instances. Its guard counts were reasonably low: on average $\sim18\%$ of the polygon's vertices were chosen as guards. In terms of optimality, greedy used about 1.17 times the minimum number of guards on average (see Table~\ref{tab:small_results_summary}), with the worst-case being 2.0 times the optimum in our dataset. This confirms that the greedy approach is a strong baseline for these instances, providing full coverage with relatively small redundancy in guards.
    \item \textbf{Reinforcement Learning:} The RL approach, using a pointer network trained with REINFORCE, achieved high coverage on many instances but not all. On average, the coverage was around 82--90\% (depending on the specific trained model checkpoint). The best model (in terms of coverage) we obtained covered about 89.9\% of the area on average for the validation set. However, there was variance: some instances were fully covered by the RL solution (coverage = 1.0), while others were only partially covered (in a few outliers, $<10\%$ coverage, though those were rare failures). The RL model's guard usage tended to be higher than greedy in these cases. On average, the RL solutions used about 1.7 times the optimal number of guards (Table~\ref{tab:small_results_summary}). This is worse than the greedy's 1.17, indicating the RL learned policy does not yet match the efficiency of the classical heuristic. We also note a higher standard deviation in the RL method's relative guard count (occasionally the RL placed far too many guards due to idiosyncratic learned behavior for certain polygons, or conversely placed too few and left large areas uncovered).
    \item \textbf{Supervised Learning:} The supervised approach on these small instances produced near-full coverage in most cases. As shown in Table~\ref{tab:results_small}, for instance size ranges 10--59, 60--99, and 100--199 vertices, the mean coverage of the SL method's solutions was 0.95, 0.93, and 0.87 respectively (averaging to about 0.89 overall). Many instances thus had almost complete coverage by the guards suggested by the network. However, this came at a great cost in terms of number of guards: the average guard set size was very large relative to optimum. In the smallest category (10--59 vertices), the SL method used on average 6.26 times the optimal number of guards; this ratio worsened for larger polygons in the same group (e.g., 19.89 times optimal for 100-199 vertices). The supervised model often selected a vast majority of vertices as guards, essentially erring on the side of covering everything with redundancy. While this ensures high coverage, it is highly suboptimal in guard count. The distribution of the coverage and guard count ratios for these instances is depicted in Figure~\ref{fig:coverage_ratio_boxplot}(a). We see that coverage is generally high (median near 1.0) for the SL method on small instances, but the ratio to optimum is extremely high (median well above 1, with large variance).
\end{itemize}

\begin{table}[!h]
\centering
\scriptsize
\begin{tabular}{lccc}
\toprule
\textbf{Method} & \textbf{Average Coverage} & \textbf{Average \#Guards} & \textbf{Std. Dev. (coverage)} \\
\midrule
Greedy & 1.000 & 1.17 $\times$ optimum & 0.0 \\
REINFORCE & 0.82--0.90 & $\sim$1.7 $\times$ optimum & 0.20--0.28 \\
Supervised & 0.89 & $\sim$10.0 $\times$ optimum & 0.07 \\
\bottomrule
\end{tabular}
\caption{Summary of performance on small (up to 200-vertex) gallery instances. Greedy achieves perfect coverage with solutions close to optimal guard count. The RL approach has high but not complete coverage and uses more guards than optimal on average. The supervised approach achieves near-high coverage but uses an order of magnitude more guards than optimal. (For RL, a range is given due to variability between different training runs/models; we report the range covering the best models we obtained.)}
\label{tab:small_results_summary}
\end{table}

\begin{figure}[h]
  \centering
  \includegraphics[width=0.8\textwidth]{gfx//method_comparison.png}
  \caption{Your figure caption}
  \label{fig:your_label}
\end{figure}


\vspace{1ex}

\noindent \textbf{Example Solution Illustration:} Figure~\ref{fig:demo_solutions} provides a qualitative comparison on one gallery instance. The optimal solution (a) uses 17 guards to cover the polygon. The RL generator model produced several candidate solutions; one such candidate (b) used 16 guards and achieved 94\% coverage (missing a small area), which corresponds to 0.94 of the optimal guard count (since it used slightly fewer guards than the optimum) and 94\% area coverage. Other candidates (c--f) generated by the RL model demonstrate the trade-offs: some used far fewer guards (down to 10 guards in (f), which is 58\% of optimum) but consequently had much lower coverage (82\%), while others used 15 guards with moderate coverage (~89--93\%). In this example, the CovNet evaluator identified (b) as the best candidate, which is indeed very close to the optimal solution. This highlights that, given a good evaluator, the RL method can generate near-optimal solutions among its outputs, even if its single-run performance may vary.

\begin{figure}[!h]
    \centering
    \subfigure[Optimal (17 guards)]{\includegraphics[scale=0.4]{gfx/demo/opt.png}}
    \hfill
    \subfigure[CovNet pick: 16 guards, cov. 0.94]{\includegraphics[scale=0.4]{gfx/demo/covnet_predicted.png}}
    \hfill
    \subfigure[Candidate 1: 10 guards, cov. 0.72]{\includegraphics[scale=0.4]{gfx/demo/cand1.png}}
    \hfill
    \subfigure[Candidate 2: 15 guards, cov. 0.89]{\includegraphics[scale=0.4]{gfx/demo/cand2.png}} \\
    \subfigure[Candidate 3: 15 guards, cov. 0.93]{\includegraphics[scale=0.4]{gfx/demo/cand3.png}}
    \hfill
    \subfigure[Candidate 4: 10 guards, cov. 0.82]{\includegraphics[scale=0.4]{gfx/demo/cand4.png}}
    \caption{An example gallery instance demonstrating solutions by different methods. (a) shows the optimal solution with 17 guards (each guard's visibility region is implicitly indicated by coverage of the polygon). (b) shows a high-quality solution generated by the RL model (16 guards) which was selected by the CovNet evaluator as having the best predicted coverage (94\% coverage, very close to full). (c)--(f) show other candidate solutions generated by the RL model for the same instance, illustrating the range of outcomes (each subcaption lists the number of guards and achieved coverage). The supervised approach, in contrast, would typically place a guard at nearly every vertex for this instance (not shown here), achieving full coverage but with far more than 17 guards.}
    \label{fig:demo_solutions}
\end{figure}

\subsection{Generalization to Larger Instances}
One concern for learning-based approaches is how well they generalize beyond the conditions on which they were trained. We tested the methods on a set of larger gallery polygons (200--600 vertices) that were outside the training range of the RL and SL models. The greedy algorithm, being an instance-specific heuristic, continued to perform well on these larger instances: it still achieved 100\% coverage on all cases, and the relative guard count remained similar (in fact, for larger polygons, greedy often performed within 1.1--1.2 times optimum on average, with some worst-cases around 1.5 times optimum, possibly because larger polygons offer more opportunities for each guard to cover area efficiently).

The RL model in its trained form did not directly generalize to significantly larger polygons if they were far beyond the training distribution (e.g., a model trained on up to 200-vertex polygons struggled on 500-vertex ones, often failing to cover large portions as it had not learned the strategies for much bigger spaces). In principle, one could train a separate RL model on larger instances or use curriculum learning to extend the range, but that was outside our current scope. Thus, we mainly rely on the supervised model for a learning approach on the larger instances.

The supervised CovNet-based method was applied to the large instances. Here, the results were notably poorer in coverage. As summarized in Table~\ref{tab:results_large}, the mean coverage achieved by the supervised model on polygons with 400-500 vertices was around 0.59 (59\%), dropping from ~0.86-0.77 for moderately larger polygons (200-399 vertices). The performance gradually degrades as instance size grows beyond what the model saw in training. This indicates that the model did not fully capture a scale-invariant strategy for guard placement; instead, it likely overfit to patterns in smaller layouts. The number of guards it placed, however, remained extremely high relative to optimum (averaging 23-27 times the optimal guard count in these large cases). Essentially, the model was placing guards almost everywhere it could, yet still missing some spots because it had not learned how to systematically ensure full coverage in larger spaces.

Table~\ref{tab:results_large} details these outcomes by binned polygon sizes. Figure~\ref{fig:coverage_ratio_boxplot}(b) illustrates the distribution of coverage and optimal ratio for the supervised method on these larger instances. In the box plot, we see a wide spread in coverage (some large instances were relatively well-covered, others very poorly) and consistently huge values for the guard count ratio (with minima still far above 1.0).

\begin{table}[!h]
\centering
\small 
    \begin{tabular}{lrrrrr}
        \toprule
        \textbf{Polygon size range} &  \textbf{Instance count} &  \multicolumn{2}{c}{\textbf{Coverage}} &  \multicolumn{2}{c}{\textbf{Ratio to Optimum}} \\
        & &  mean &  std  &  mean  &  std \\
        \midrule
         10--59     &     149 &           0.95 &          0.08 &        6.26 &       2.43 \\
         60--99     &     219 &           0.93 &          0.06 &       11.75 &       2.53 \\
         100--199   &     853 &           0.87 &          0.07 &       19.89 &       4.19 \\
        \midrule
        \textit{(200--299)}    &     373 &           0.86 &          0.09 &       23.80 &       5.07 \\
        \textit{(300--399)}    &     202 &           0.77 &          0.17 &       25.85 &       5.67 \\
        \textit{(400--499)}    &     180 &           0.64 &          0.20 &       27.04 &       5.03 \\
        \textit{(500--599)}    &     129 &           0.59 &          0.23 &       27.59 &       5.29 \\
        \bottomrule
    \end{tabular}
\caption{Performance of the \textbf{supervised learning (CovNet) method} on two groups of instances: The top rows show results for smaller instances (10--199 vertices, totaling 1221 instances), and the bottom rows (in italics) show results for larger instances (200--599 vertices, 884 instances). \emph{Coverage} is the fraction of area covered by the guards chosen by the model, and \emph{Ratio to Optimum} is the number of guards used divided by the optimal number of guards for that instance. We observe high coverage and moderately low variance on smaller instances, at the expense of large redundant guard sets. On larger instances, coverage drops and the guard count ratio further increases, indicating poor scalability of the supervised model's approach.}
\label{tab:results_small}\label{tab:results_large}
\end{table}

\begin{figure}[!h]
    \centering
    \subfigure[Instances of sizes seen during training]{\includegraphics[scale=0.5]{gfx/testbed_1.png}}
    \hfill
    \subfigure[Larger instances (unseen sizes)]{\includegraphics[scale=0.5]{gfx/testbed_2.png}}
    \caption{Distribution of coverage and guard count ratio (to optimal) for the \textbf{supervised learning method}'s solutions on two test sets: (a) smaller instances within the training distribution (10--199 vertices), and (b) larger instances beyond what the model was trained on (200--599 vertices). Each box plot summarizes the spread of coverage achieved (left in each sub-figure) and the ratio of the solution's guard count to the optimal guard count (right in each sub-figure). In (a), the model usually achieves high coverage (many instances at or near 1.0 coverage) but the guard count ratio is very high (median in double digits). In (b), the coverage median drops (significantly lower, indicating many large instances are only partially covered), while the guard count ratio remains extremely high.}
    \label{fig:coverage_ratio_boxplot}
\end{figure}

These generalization results highlight a key challenge: \textbf{the learned models (both RL and SL) are limited by the distribution of their training data and objectives, whereas the greedy algorithm is not}. The greedy method inherently scales to any polygon (within computational limits) and always guarantees full coverage, though it has no guarantee on optimality. The RL method can, in principle, be retrained or fine-tuned for larger instances, but that requires additional training effort and does not guarantee success if the complexity grows. The supervised method as implemented clearly did not learn a general strategy applicable to much larger polygons, indicating perhaps the need for architectures that capture more global or scalable features, or training regimes that include more diverse and larger examples.

\subsection{Discussion of Trade-offs}
From our comparative study, several insights emerge:
\begin{itemize}
    \item \textbf{Solution Quality vs Optimal:} The greedy algorithm provided solutions closest to optimal in guard count while guaranteeing full coverage. The RL approach provided somewhat worse coverage and guard efficiency; its advantage would only be justified if it could significantly improve runtime or scale better, but in our experiments the RL model did not clearly outperform the greedy in quality. The supervised approach produced the worst guard efficiency, essentially trading optimality for coverage (and even then failing to cover fully on large cases).
    \item \textbf{Computational Considerations:} The greedy algorithm, being a hand-crafted heuristic, was moderately fast on our instance sizes (taking seconds per instance). The RL and SL methods have a large upfront cost (training can take hours), but once trained, producing a solution is very fast (a single network forward pass taking milliseconds). This suggests that if one needed to solve a very large number of instances repeatedly, a learned approach might become more appealing. In scenarios where a model can be trained offline and then deployed in real-time (e.g., a robot that needs to quickly decide guard or camera placements on the fly), the learned methods offer an advantage of speed at runtime.
    \item \textbf{Robustness and Guarantees:} The greedy method has a predictable performance: it will always cover the area completely. The RL and SL methods lack guarantees -- they might fail to cover some parts of the polygon and there's no easy way to know without computing coverage. In safety-critical applications, this is a major drawback. One could combine approaches (e.g., run the learned method and then use a quick greedy top-up to cover any uncovered spots, thereby ensuring coverage while hopefully using fewer guards than a pure greedy approach). Such hybrid strategies could marry the strengths of both sides.
    \item \textbf{Learning Problem Difficulty:} The disparity in performance between the learned models and the greedy algorithm underscores the difficulty of the AGP as a learning problem. Unlike simpler combinatorial problems (like shortest paths or even traveling salesman for moderate sizes) where learning approaches have come closer to classical methods, the AGP involves a complex continuous space and combinatorial coverage aspect that is hard to capture with limited training. The supervised model effectively learned to cover the polygon but not to do so efficiently, indicating that the optimal coverage problem (cover everything with as few guards as possible) might require more sophisticated training signals or model architectures to learn effectively. The RL model, which had an explicit reward balancing coverage and count, still struggled to find the true optimal balance, likely due to the huge search space of possible guard sets.
    \item \textbf{Potential Improvements:} There are ways to improve the learning-based methods. For RL, one could incorporate more advanced techniques like curriculum learning (starting with smaller polygons then gradually increasing size, which we partially did), use of baseline or critic networks to reduce variance in training, or using more structured action representations. For supervised learning, using graph neural networks or transformers that operate on the polygon geometry rather than a fixed grid might improve generalization to larger or different shapes, as they could be more size-invariant. Additionally, explicitly training to minimize guard count (perhaps by generating not just any covering set but minimal covers) could push the model to be more efficient. These are avenues for future research.
\end{itemize}

\section{Conclusion}
\label{sec:conclusion}
We presented a comparative study of three methods for solving the Art Gallery Problem: a greedy algorithm, a reinforcement learning approach, and a supervised learning approach. Our evaluation on a common set of test polygons revealed that the classical greedy heuristic still outperforms the learning-based methods in terms of guaranteed coverage and efficiency of the solution (guard count). The reinforcement learning method, while promising and able to learn non-trivial strategies, did not consistently achieve full coverage and tended to use more guards than optimal. The supervised learning method, as implemented via a convolutional network heuristic, achieved high coverage on smaller instances but was extremely inefficient in terms of guard usage and failed to generalize well to larger instances.

These results suggest that, for the Art Gallery Problem, current learning-based approaches have not yet matched the reliability and quality of a well-designed algorithmic heuristic for the instances we tested. Nonetheless, the machine learning approaches offer benefits in terms of fast deployment and the potential to improve with more training data or better architectures. They also open the door to tackling variants of the problem where classical methods struggle (for example, when the environment has uncertainty, or the guard visibility has complex constraints, a learned approach might adapt more readily).

In conclusion, greedy and other algorithmic methods remain the go-to solution for obtaining high-quality results in AGP, but reinforcement learning and supervised models are an intriguing alternative that with further development could become competitive. Future work should explore hybrid methods (combining learned components with algorithmic guarantees) and investigate learning strategies that focus on the combinatorial optimality (not just feasibility) of solutions. The Art Gallery Problem continues to be a challenging testbed for both algorithms and learning, and insights from each domain can inform the other.

\bibliographystyle{plain}
\bibliography{paper}
\end{document}
